\section{Wick's theorem and the Feynman rules}\label{app:D.11}

Every diagram of \cref{sec:4.8} is one term of an expansion of the Green's function in powers of the interaction. Each term is an average of a long string of free creation and annihilation operators, and Wick's theorem reduces that average to a sum of products of free Green's functions, one product for each way of pairing the operators. This section derives the expansion in imaginary time, where it is exact at any temperature, and at $T = 0$ in real time, where it rests on the adiabatic theorem of Gell-Mann and Low~\cite{gellmann1951bound}. It then proves Wick's theorem~\cite{wick1950evaluation} by Gaudin's method~\cite{gaudin1960demonstration}, shows that the disconnected terms cancel, and derives the Feynman rules for the Coulomb interaction and for the exchange of a boson, with every sign and every combinatorial factor. The conventions for $\GF$ are those of \cref{sec:4.7}: a bare $\GF$ with a real argument is the time-ordered function, and with an imaginary argument the Matsubara function. The treatment follows Refs.~\cite[Chaps.~3 and 7]{fetter1971quantum} and \cite{bruus2004many}.

\subsection*{The expansion in imaginary time}

Split $H_\mu = H_0 + H_{\mathrm{int}}$, with $H_0 = \sum_{\vect{k}}\xi_{\vect{k}}c^\dagger_{\vect{k}}c_{\vect{k}}$, plus $\sum_{\vect{Q}}\omega_{\vect{Q}}d^\dagger_{\vect{Q}}d_{\vect{Q}}$ when bosons are present. Operators in the interaction picture evolve with $H_0$ alone, $O_{\mathrm{I}}(\tau) = \ee^{\tau H_0}O\,\ee^{-\tau H_0}$, so that $c_{\vect{k},\mathrm{I}}(\tau) = \ee^{-\xi_{\vect{k}}\tau}c_{\vect{k}}$ and $c^\dagger_{\vect{k},\mathrm{I}}(\tau) = \ee^{\xi_{\vect{k}}\tau}c^\dagger_{\vect{k}}$. Define $U_{\mathrm{I}}(\tau,0) = \ee^{\tau H_0}\ee^{-\tau H_\mu}$. Its derivative is $\ee^{\tau H_0}(H_0 - H_\mu)\ee^{-\tau H_\mu} = -H_{\mathrm{int},\mathrm{I}}(\tau)\,U_{\mathrm{I}}(\tau,0)$, which has the form of \cref{eq:2.4-ipicture} with $\ii\,\dd t$ replaced by $-\dd\tau$. Integrating and iterating as in \cref{eq:2.4-dyson},
\begin{equation*}
  \begin{aligned}
  U_{\mathrm{I}}(\tau,0) &= \sum_{n=0}^\infty(-1)^n\int_0^\tau\dd\tau_1\int_0^{\tau_1}\dd\tau_2\cdots\int_0^{\tau_{n-1}}\dd\tau_n\,H_{\mathrm{int},\mathrm{I}}(\tau_1)\cdots H_{\mathrm{int},\mathrm{I}}(\tau_n) \\
  &= \sum_{n=0}^\infty\frac{(-1)^n}{n!}\int_0^\tau\dd\tau_1\cdots\dd\tau_n\,\mathrm{T}_\tau\bigl[H_{\mathrm{int},\mathrm{I}}(\tau_1)\cdots H_{\mathrm{int},\mathrm{I}}(\tau_n)\bigr] .
  \end{aligned}
\end{equation*}
In the second form each integral runs over the whole interval. The $n!$ orderings of the times give the same ordered product, since $\mathrm{T}_\tau$ restores the order and $H_{\mathrm{int}}$ contains fermion operators in pairs, so that no sign arises. The same steps give $U_{\mathrm{I}}(\tau,\tau') = \ee^{\tau H_0}\ee^{-(\tau - \tau')H_\mu}\ee^{-\tau'H_0}$ as the ordered exponential over $[\tau', \tau]$.

The Boltzmann factor and the Heisenberg operators of \cref{sec:4.7} are both expressed through $U_{\mathrm{I}}$. The Boltzmann factor is $\ee^{-\beta H_\mu} = \ee^{-\beta H_0}U_{\mathrm{I}}(\beta,0)$, and the full imaginary-time operator is
\begin{equation*}
  c_{\vect{k}}(\tau) = \ee^{\tau H_\mu}c_{\vect{k}}\ee^{-\tau H_\mu} = U_{\mathrm{I}}(\tau,0)^{-1}\,c_{\vect{k},\mathrm{I}}(\tau)\,U_{\mathrm{I}}(\tau,0) .
\end{equation*} For $0 < \tau < \beta$, inserting both into \cref{eq:4.7-Gtau} and using $U_{\mathrm{I}}(\beta,0)U_{\mathrm{I}}(\tau,0)^{-1} = U_{\mathrm{I}}(\beta,\tau)$,
\begin{equation*}
  \Tr\bigl[\ee^{-\beta H_\mu}c_{\vect{k}}(\tau)c^\dagger_{\vect{k}}\bigr] = \Tr\bigl[\ee^{-\beta H_0}\,U_{\mathrm{I}}(\beta,\tau)\,c_{\vect{k},\mathrm{I}}(\tau)\,U_{\mathrm{I}}(\tau,0)\,c^\dagger_{\vect{k},\mathrm{I}}(0)\bigr] .
\end{equation*}
Every factor on the right stands to the left of every factor with an earlier imaginary time, so the product is the time-ordered product of $U_{\mathrm{I}}(\beta,0)$ with the two operators. The case $\tau < 0$ works the same way, and the partition function is $\mathcal{Z} = \mathcal{Z}_0\langle U_{\mathrm{I}}(\beta,0)\rangle_0$, where $\langle\cdots\rangle_0 = \Tr[\ee^{-\beta H_0}\cdots]/\mathcal{Z}_0$ is the average in the free grand-canonical state. Hence
\begin{equation}
  \GF(\vect{k},\tau) = -\frac{\bigl\langle\mathrm{T}_\tau\,U_{\mathrm{I}}(\beta,0)\,c_{\vect{k},\mathrm{I}}(\tau)\,c^\dagger_{\vect{k},\mathrm{I}}(0)\bigr\rangle_0}{\bigl\langle U_{\mathrm{I}}(\beta,0)\bigr\rangle_0} .
  \label{eq:D.11-Gexp}
\end{equation}
Nothing has been assumed beyond the convergence of the series. The density response \cref{eq:4.7-Rtau} has the same expansion with $n_e(\vect{q})^\dagger$ and $n_e(\vect{q})$ in place of $c_{\vect{k}}$ and $c^\dagger_{\vect{k}}$, and a factor $1/\Vol$.

Inside $H_{\mathrm{int},\mathrm{I}}(\tau)$ creation and annihilation operators share one time, and the ordering symbol needs a rule for them. We keep them in the order in which they are written in $H_{\mathrm{int}}$, creation operators to the left; equivalently, every creation operator in $H_{\mathrm{int}}$ is taken at $\tau + 0^+$.

\subsection*{The expansion at zero temperature}

At $T = 0$ the time-ordered function \cref{eq:4.7-GT} is an average in the interacting ground state $|0\rangle$, which is not known. Switch the interaction on and off slowly, $H_{\mathrm{int}}\ee^{-\eta_{\mathrm{b}}|t|}$ with $\eta_{\mathrm{b}} \to 0^+$, and start from the ground state $|\Phi_0\rangle$ of $H_0$, the Fermi sea. The theorem of Gell-Mann and Low~\cite{gellmann1951bound} states that $U_{\mathrm{I}}(0,-\infty)|\Phi_0\rangle/\langle\Phi_0|U_{\mathrm{I}}(0,-\infty)|\Phi_0\rangle$, with $U_{\mathrm{I}}$ the real-time operator of \cref{eq:2.4-dyson} with $H' = H_{\mathrm{int}}$, tends to an eigenstate of $H_\mu$. It is the ground state when that state is nondegenerate and is reached from $|\Phi_0\rangle$ without a level crossing as the interaction is switched on, which we assume. Switching off toward $t = +\infty$ leads to the same state, up to a phase, so $\langle\Phi_0|U_{\mathrm{I}}(\infty,0)$ is proportional to $\langle0|$. With $c_{\vect{k}}(t) = U_{\mathrm{I}}(t,0)^{-1}c_{\vect{k},\mathrm{I}}(t)\,U_{\mathrm{I}}(t,0)$ the steps of the previous paragraph go through with real times, and the unknown phases and normalizations cancel in the ratio:
\begin{equation}
  \GF(\vect{k},t) = -\ii\,\frac{\langle\Phi_0|\mathrm{T}\,U_{\mathrm{I}}(\infty,-\infty)\,c_{\vect{k},\mathrm{I}}(t)\,c^\dagger_{\vect{k},\mathrm{I}}(0)|\Phi_0\rangle}{\langle\Phi_0|U_{\mathrm{I}}(\infty,-\infty)|\Phi_0\rangle} ,
  \label{eq:D.11-GT}
\end{equation}
with $U_{\mathrm{I}}(\infty,-\infty) = \sum_n(-\ii)^n(n!)^{-1}\int\dd t_1\cdots\dd t_n\,\mathrm{T}[H_{\mathrm{int},\mathrm{I}}(t_1)\cdots H_{\mathrm{int},\mathrm{I}}(t_n)]$. At $T > 0$ no expansion of this form exists for real times, since a thermal state is not reached adiabatically from a single free state. A real-time expansion at finite temperature needs a time contour that runs forward and back, and the imaginary-time form \cref{eq:D.11-Gexp} replaces it here (\cref{sec:4.7}).

\subsection*{Wick's theorem}

Let $A_1, \dots, A_{2n}$ be operators of the free problem, each a creation or an annihilation operator of one mode, multiplied by a number. In imaginary time $c_{\vect{k},\mathrm{I}}(\tau)$ is of this kind, and so is $c_{\vect{k},\mathrm{I}}(t)$ in real time. For fermions, the claim is
\begin{equation}
  \langle A_1A_2\cdots A_{2n}\rangle_0 = \sum_{j=2}^{2n}(-1)^j\,\langle A_1A_j\rangle_0\,\langle A_2\cdots\widehat{A_j}\cdots A_{2n}\rangle_0 ,
  \label{eq:D.11-wick}
\end{equation}
where the hat removes $A_j$ from the string. Applied repeatedly, it writes the average as a sum over all ways of splitting the $2n$ operators into pairs. Each pairing contributes the product of its pair averages, with the pair kept in its original order, and the sign of the permutation that brings the members of every pair next to each other. An average of an odd number of operators vanishes, and so does a pair average unless it holds one creation and one annihilation operator of the same mode.

The proof uses only the anticommutators and the cyclic property of the trace~\cite{gaudin1960demonstration}. Move $A_1$ to the right end, one step at a time. Each step past $A_j$ replaces $A_1A_j$ by $\{A_1, A_j\} - A_jA_1$, and the anticommutator is a number. After $2n - 1$ steps,
\begin{equation*}
  A_1A_2\cdots A_{2n} = \sum_{j=2}^{2n}(-1)^j\{A_1, A_j\}\,A_2\cdots\widehat{A_j}\cdots A_{2n} - A_2\cdots A_{2n}A_1 .
\end{equation*}
In the average of the last term, the trace moves $A_1$ back to the front, past the Boltzmann factor: $\Tr[\ee^{-\beta H_0}A_2\cdots A_{2n}A_1] = \Tr[A_1\ee^{-\beta H_0}A_2\cdots A_{2n}]$. If $A_1$ annihilates a fermion of energy $\xi$, then $A_1\ee^{-\beta H_0} = \ee^{-\beta\xi}\ee^{-\beta H_0}A_1$, because removing the fermion lowers $H_0$ by $\xi$; if it creates one, the factor is $\ee^{+\beta\xi}$. Writing $\ee^{-\beta\xi_1}$ for either case, with $\xi_1 = \xi$ for an annihilation and $\xi_1 = -\xi$ for a creation operator, the last term averages to $-\ee^{-\beta\xi_1}\langle A_1\cdots A_{2n}\rangle_0$, and
\begin{equation*}
  \langle A_1\cdots A_{2n}\rangle_0 = \sum_{j=2}^{2n}(-1)^j\,\frac{\{A_1, A_j\}}{1 + \ee^{-\beta\xi_1}}\,\langle A_2\cdots\widehat{A_j}\cdots A_{2n}\rangle_0 .
\end{equation*}
For $n = 1$ the same equation reads $\langle A_1A_j\rangle_0 = \{A_1, A_j\}/(1 + \ee^{-\beta\xi_1})$; for example, $\langle c_{\vect{k}}c^\dagger_{\vect{k}}\rangle_0 = 1/(1 + \ee^{-\beta\xi_{\vect{k}}}) = 1 - n_{\mathrm{F}}(\xi_{\vect{k}})$. Inserting this in the general case gives \cref{eq:D.11-wick}. For bosons every anticommutator becomes a commutator, no signs appear, and the denominator is $1 - \ee^{-\beta\xi_1}$, which gives $\langle d_{\vect{Q}}d^\dagger_{\vect{Q}}\rangle_0 = 1 + n_{\mathrm{B}}(\omega_{\vect{Q}})$. Fermion and boson operators commute and $H_0$ is a sum of a fermion and a boson part, and a mixed average factorizes into a fermion average and a boson average. Linear combinations of single-mode operators, such as $d_{\vect{Q}} + d^\dagger_{-\vect{Q}}$, are covered by linearity.

The expansions need the theorem for time-ordered strings and at $T = 0$. Inside $\mathrm{T}_\tau$ the operators are first put in time order, at the cost of the sign of that permutation, and \cref{eq:D.11-wick} is applied to the ordered string. Each pair then appears in time order, as $\langle\mathrm{T}_\tau A_iA_j\rangle_0$, and the sign of the time ordering combines with the sign of the pairing in the ordered string into the sign of the same pairing in the original string. Hence the time-ordered average is the sum over pairings of products of time-ordered pair averages, the contractions, with the sign of the pairing as written. The free contractions are, by \cref{eq:4.7-Gtau} and \cref{eq:4.7-GT},
\begin{equation}
  \langle\mathrm{T}_\tau\,c_{\vect{k},\mathrm{I}}(\tau)\,c^\dagger_{\vect{k},\mathrm{I}}(\tau')\rangle_0 = -\GF_0(\vect{k},\tau - \tau'), \qquad
  \langle\Phi_0|\mathrm{T}\,c_{\vect{k},\mathrm{I}}(t)\,c^\dagger_{\vect{k},\mathrm{I}}(t')|\Phi_0\rangle = \ii\,\GF_0(\vect{k},t - t') .
  \label{eq:D.11-contract}
\end{equation}
At $T = 0$, the ground state $|\Phi_0\rangle$ of $H_0$ is the limit $\beta \to \infty$ of the free grand-canonical state, every average in \cref{eq:D.11-wick} has a limit, and the theorem holds for $\langle\Phi_0|\cdots|\Phi_0\rangle$, in real time as in imaginary time.

\subsection*{Disconnected terms cancel}

At order $n$, each pairing in the numerator of \cref{eq:D.11-Gexp} joins the two external operators to some set of $m$ of the $n$ factors $H_{\mathrm{int}}$, the connected part, and joins the remaining $n - m$ factors only among themselves, the vacuum part. The value of the pairing is the product of the two parts. Since every $H_{\mathrm{int}}$ is even in fermion operators, its position under $\mathrm{T}_\tau$ costs no sign, and the sign of the pairing is the product of the signs of the two parts. There are $n!/[m!\,(n - m)!]$ ways to choose which of the $n$ identical, integrated factors form the connected part, and each gives the same value. With $X_m$ the sum of connected pairings at order $m$, including the external operators, and $Y_l$ the sum of all pairings of $l$ factors $H_{\mathrm{int}}$ alone, the numerator is
\begin{equation*}
  \sum_n\frac{(-1)^n}{n!}\sum_{m=0}^n\frac{n!}{m!\,(n - m)!}X_m\,Y_{n-m}
  = \Bigl[\sum_m\frac{(-1)^m}{m!}X_m\Bigr]\Bigl[\sum_l\frac{(-1)^l}{l!}Y_l\Bigr] .
\end{equation*}
The second factor is $\langle U_{\mathrm{I}}(\beta,0)\rangle_0$, which cancels the denominator. Hence $\GF$ is $-1$ times the sum of the connected pairings alone, the linked-cluster theorem. The same argument applies to \cref{eq:D.11-GT} with $(-\ii)^n$ in place of $(-1)^n$.

\subsection*{Feynman rules for the Coulomb interaction}

Take the electron gas of \cref{eq:2.3-jellium}, whose interaction is
\begin{equation*}
  H_{\mathrm{int}} = \frac{1}{2\Vol}\sum_{\vect{k}\ne0}v_{\mathrm{C}}(k)\sum_{\vect{p}\vect{p}'ss'}c^\dagger_{\vect{p}+\vect{k},s}c^\dagger_{\vect{p}'-\vect{k},s'}c_{\vect{p}'s'}c_{\vect{p}s} .
\end{equation*} Under $\mathrm{T}_\tau$, with the creation operators at $\tau + 0^+$, the order of the four operators is fixed by the ordering symbol, and they may be regrouped with the sign of the permutation. Moving $c_{\vect{p}s}$ two places to the left costs no sign, and each factor is the product of two bilinears,
\begin{equation*}
  H_{\mathrm{int},\mathrm{I}}(\tau) = \frac{1}{2\Vol}\sum_{\vect{k}\ne0}v_{\mathrm{C}}(k)\sum_{\vect{p}\vect{p}'ss'}\bigl[c^\dagger_{\vect{p}+\vect{k},s}c_{\vect{p}s}\bigr]\bigl[c^\dagger_{\vect{p}'-\vect{k},s'}c_{\vect{p}'s'}\bigr](\tau) \qquad\text{(inside $\mathrm{T}_\tau$)} ,
\end{equation*}
each bilinear a density component. We draw each factor as a dashed interaction line carrying the momentum $\vect{k}$, whose two ends are the bilinears, and each contraction \cref{eq:D.11-contract} as a directed solid line from a creation operator to the annihilation operator it is paired with. At order $n$ there are $2n$ bilinears and the two external operators, which makes $2n + 1$ contractions.

\emph{The sign.} Follow a fermion line from the external $c_{\vect{k}}(\tau)$ through the bilinears it visits to $c^\dagger_{\vect{k}}(0)$. Arrange the string so that the bilinears it visits stand in the order of the line, $c(\tau)\,[c^\dagger_1c_1]\,[c^\dagger_2c_2]\cdots[c^\dagger_mc_m]\,c^\dagger(0)$; moving whole bilinears costs no sign. The contractions $c(\tau)c^\dagger_1$, $c_1c^\dagger_2$, \dots, $c_mc^\dagger(0)$ are then all between neighbors, the pairing sign is $+1$, and each contraction is $-\GF_0$, which makes $(-1)^{m+1}$ for the $m + 1$ propagators of the open line. A closed loop through $m$ bilinears, $[c^\dagger_1c_1]\cdots[c^\dagger_mc_m]$, is paired as $c_1c^\dagger_2$, \dots, $c_{m-1}c^\dagger_m$ and finally $c_mc^\dagger_1$. All pairs are neighbors except the last. Bringing $c^\dagger_1$ from the front to the end passes $2m - 1$ operators and costs a factor $-1$, in addition to $(-1)^m$ from the $m$ propagators of the loop. With $F$ closed loops, the $2n + 1$ propagators therefore carry $(-1)^{2n+1}(-1)^F = -(-1)^F$. The definition \cref{eq:D.11-Gexp} contributes $-1$ and the expansion $(-1)^n$. The total sign of a diagram of order $n$ is $(-1)^n(-1)^F$.

\emph{The factors $1/n!$ and $1/2^n$.} Relabeling the $n$ integrated interaction factors among themselves, and exchanging the two ends of any interaction line, which is a relabeling $\vect{k} \to -\vect{k}$, $(\vect{p}s) \leftrightarrow (\vect{p}'s')$ under which $v_{\mathrm{C}}$ is unchanged, maps one pairing onto another with the same value. For a connected diagram with two external points the $n!\,2^n$ relabelings give $n!\,2^n$ different pairings, since the external points fix every line of the diagram. Summing over pairings is then the same as summing over topologically distinct diagrams with weight $n!\,2^n$, which cancels the $1/n!$ of the expansion and the $\tfrac12$ in each $H_{\mathrm{int}}$.

\emph{Frequencies.} Each propagator is written as $\GF_0(\vect{p},\tau_a - \tau_b) = T\sum_j\ee^{-\ii\omega^{\mathrm{F}}_j(\tau_a - \tau_b)}\GF_0(\vect{p},\ii\omega^{\mathrm{F}}_j)$ by \cref{eq:4.7-Gmats}. The four ends of an interaction line share one time $\tau_v$. The integral $\int_0^\beta\dd\tau_v$ of the product of their exponentials is $\beta$ times a Kronecker delta that conserves frequency along the line, which therefore carries the bosonic difference of the two fermionic frequencies at either end. The transform $\int_0^\beta\dd\tau\,\ee^{\ii\omega^{\mathrm{F}}_j\tau}$ of the external time supplies one more $\beta$. Of the $2n + 1$ factors $T$, $n + 1$ are cancelled by these factors $\beta$, and $n$ remain, one for each independent internal frequency. Momenta are conserved at each end of a line in the same way, and the $n$ factors $1/\Vol$ go with the $n$ independent internal momentum sums.

Collecting, the Matsubara Green's function $\GF(\vect{k},\ii\omega^{\mathrm{F}}_j)$ at order $n$ is the sum over connected, topologically distinct diagrams with $n$ interaction lines, each diagram evaluated with the rules listed in \cref{sec:4.8}. The interaction line $-v_{\mathrm{C}}(k)/\Vol$ is independent of $\omega^{\mathrm{B}}$, and the factor $\ee^{\ii\omega^{\mathrm{F}}0^+}$ of a solid line that begins and ends on the same interaction line comes from the creation operator at $\tau + 0^+$.
The self-energy $\Sigma(\vect{k},\ii\omega^{\mathrm{F}}_j)$ is the sum of the diagrams that cannot be cut into two by cutting one solid line, with the two external propagators removed. The Dyson equation of \cref{sec:4.7} then sums all chains of them.

At first order the self-energy has a Hartree and a Fock diagram. In the Hartree diagram the interaction line carries $\vect{k} = 0$ into a closed loop; it is $(-1)(-1)\,[v_{\mathrm{C}}(0)/\Vol]\allowbreak\sum_{\vect{p}s}T\sum_j\ee^{\ii\omega^{\mathrm{F}}_j0^+}\GF_0(\vect{p},\ii\omega^{\mathrm{F}}_j) = v_{\mathrm{C}}(0)N_e/\Vol$ by \cref{eq:4.7-occmats}, the electrostatic potential of the uniform density, and it is absent in jellium, where the background removes $\vect{k} = 0$. In the Fock diagram the line carries $\vect{k} - \vect{p}$ and there is no loop:
\begin{equation}
  \Sigma_{\mathrm{x}}(\vect{k}) = -\frac{1}{\Vol}\sum_{\vect{p}}v_{\mathrm{C}}(|\vect{k} - \vect{p}|)\,T\sum_j\ee^{\ii\omega^{\mathrm{F}}_j0^+}\GF_0(\vect{p},\ii\omega^{\mathrm{F}}_j) = -\frac{1}{\Vol}\sum_{\vect{p}}v_{\mathrm{C}}(|\vect{k} - \vect{p}|)\,n_{\mathrm{F}}(\xi_{\vect{p}}) ,
  \label{eq:D.11-fock}
\end{equation}
with $\vect{p}$ of the same spin as $\vect{k}$. At $T = 0$ this is the exchange self-energy of \cref{eq:D.4-gas}, and both signs agree with \cref{app:D.4}: the Hartree term repels and exchange attracts. The rules are thus checked at first order against a calculation that used no diagrams.

\subsection*{Exchange of a boson}

Let the electrons couple to a bosonic field with modes $\vect{Q}$ and energies $\omega_{\vect{Q}} = \omega_{-\vect{Q}}$,
\begin{equation*}
  H_{\mathrm{int}} = \frac{1}{\sqrt{\Vol}}\sum_{\vect{k}\vect{Q}s}g(\vect{Q})\,c^\dagger_{\vect{k}+\vect{Q},s}c_{\vect{k}s}\,\varphi_{\vect{Q}}, \qquad
  \varphi_{\vect{Q}} = d_{\vect{Q}} + d^\dagger_{-\vect{Q}} ,
\end{equation*}
with $g(-\vect{Q}) = g(\vect{Q})^*$, so that $H_{\mathrm{int}}$ is Hermitian, since $\varphi_{\vect{Q}}^\dagger = \varphi_{-\vect{Q}}$. For phonons $g$ becomes the vertex $g_{mn\nu}$ of \cref{sec:6.2} and $d_{\vect{Q}}$ becomes $b_{\vect{Q}\nu}$. The boson contraction defines the free boson propagator, $\mathsf{D}_0(\vect{Q},\tau) = -\langle\mathrm{T}_\tau\,\varphi_{\vect{Q},\mathrm{I}}(\tau)\,\varphi_{-\vect{Q},\mathrm{I}}(0)\rangle_0$. For $0 < \tau < \beta$, with $d_{\vect{Q},\mathrm{I}}(\tau) = \ee^{-\omega_{\vect{Q}}\tau}d_{\vect{Q}}$ and the averages of \cref{eq:2.5-nF}, it is $-[(1 + n_{\mathrm{B}})\ee^{-\omega_{\vect{Q}}\tau} + n_{\mathrm{B}}\ee^{\omega_{\vect{Q}}\tau}]$ with $n_{\mathrm{B}} = n_{\mathrm{B}}(\omega_{\vect{Q}})$. Since $\ee^{\ii\omega^{\mathrm{B}}_j\beta} = 1$, the transform is $\int_0^\beta\dd\tau\,\ee^{(\ii\omega^{\mathrm{B}}_j \mp \omega_{\vect{Q}})\tau} = (\ee^{\mp\beta\omega_{\vect{Q}}} - 1)/(\ii\omega^{\mathrm{B}}_j \mp \omega_{\vect{Q}})$, and with $(1 + n_{\mathrm{B}})(1 - \ee^{-\beta\omega_{\vect{Q}}}) = 1$ and $n_{\mathrm{B}}(\ee^{\beta\omega_{\vect{Q}}} - 1) = 1$,
\begin{equation}
  \mathsf{D}_0(\vect{Q},\ii\omega^{\mathrm{B}}_j) = \frac{1}{\ii\omega^{\mathrm{B}}_j - \omega_{\vect{Q}}} - \frac{1}{\ii\omega^{\mathrm{B}}_j + \omega_{\vect{Q}}} = \frac{2\omega_{\vect{Q}}}{(\ii\omega^{\mathrm{B}}_j)^2 - \omega_{\vect{Q}}^2} .
  \label{eq:D.11-D0}
\end{equation}
The temperature has dropped out, as it did for $\GF_0$ in \cref{sec:4.7}: thermal factors appear only after a frequency sum.

The interaction is odd in the boson operators, so only even orders $2m$ survive the boson average, and every pairing joins the $2m$ vertices into $m$ boson lines. The counting runs as for the Coulomb case. The expansion gives $(-1)^{2m}/(2m)! = 1/(2m)!$, and the $(2m)!$ relabelings of the vertices, which include the exchange of the two ends of each boson line, cancel the factorial. The fermion signs are those of $2m$ bilinears, $-(-1)^F$, and the definition of $\GF$ adds $-1$. Each boson line contributes $-\mathsf{D}_0$ from its contraction and $g(\vect{Q})g(-\vect{Q})/\Vol = |g(\vect{Q})|^2/\Vol$ from its two vertices. The total is $(-1)^m(-1)^F\prod_{\mathrm{lines}}|g|^2\mathsf{D}_0/\Vol$, which is the Coulomb result with each line replaced:
\begin{equation}
  -\frac{v_{\mathrm{C}}(k)}{\Vol} \;\longrightarrow\; -\frac{|g(\vect{Q})|^2\,\mathsf{D}_0(\vect{Q},\ii\omega^{\mathrm{B}})}{\Vol} .
  \label{eq:D.11-boson}
\end{equation}
The exchange of a boson is an interaction that depends on frequency. At zero frequency, $|g|^2\mathsf{D}_0(\vect{Q},0) = -2|g|^2/\omega_{\vect{Q}}$, an attraction. Continued to real frequencies it is $2\omega_{\vect{Q}}|g|^2/(\omega^2 - \omega_{\vect{Q}}^2)$, attractive for $|\omega| < \omega_{\vect{Q}}$ and repulsive above. An electron polarizes the field, the polarization lags behind it, and a second electron that arrives within a time $\sim 1/\omega_{\vect{Q}}$ is attracted to it. For phonons this is the attraction behind superconductivity (\cref{sec:7.2}). The first diagram of the self-energy is the Fock diagram with a boson line, which \cref{sec:4.8} evaluates (\cref{eq:4.8-sigmaboson}). Its Hartree partner carries $\vect{Q} = 0$ and is a static shift of all levels, which vanishes when $g(\vect{0}) = 0$.

\subsection*{The rules at zero temperature}

With \cref{eq:D.11-GT}, the counting is the same with three changes: the expansion gives $(-\ii)^n$, each contraction is $\ii\GF_0$ by \cref{eq:D.11-contract}, and the definition of $\GF$ carries $-\ii$. The total factor of a diagram of order $n$ is $-\ii\,(-\ii)^n\,\ii^{2n+1}(-1)^F = \ii^n(-1)^F$. Transforming with $\int\dd t\,\ee^{\ii\omega t}$ and $\GF_0(\vect{p},t) = \int\dd\omega\,(2\pi)^{-1}\ee^{-\ii\omega t}\GF_0(\vect{p},\omega)$ turns every time integral into $2\pi\delta$ of a frequency, and $n$ integrals $\int\dd\omega/2\pi$ remain. The rules for $\GF(\vect{k},\omega)$ are those of \cref{sec:4.8} with the free time-ordered propagator $\GF_0(\vect{p},\omega) = 1/(\omega - \xi_{\vect{p}} + \ii0^+\operatorname{sgn}\xi_{\vect{p}})$ of \cref{eq:4.7-GTlehmann}, each interaction line $+\ii v_{\mathrm{C}}(k)/\Vol$, and each internal frequency $\int\dd\omega/2\pi$. The Fock diagram checks the signs once more: $\ii\int\dd\omega\,(2\pi)^{-1}\ee^{\ii\omega0^+}\GF_0(\vect{p},\omega)$ is evaluated by closing the contour in the upper half plane, where the pole $\xi_{\vect{p}} + \ii0^+$ of an occupied state lies, and it equals $\ii\cdot\ii\,\Theta(-\xi_{\vect{p}}) = -\Theta(-\xi_{\vect{p}})$. With the line $v_{\mathrm{C}}/\Vol$ this is \cref{eq:D.11-fock} at $T = 0$. The time-ordered density response, $-(\ii/\Vol)\langle\mathrm{T}\,n_e(\vect{q},t)^\dagger n_e(\vect{q},0)\rangle$, carries $-\ii/\Vol$ in its definition, and its $2n + 2$ contractions give the factor $(-\ii/\Vol)(-\ii)^n\,\ii^{2n+2}(-1)^F = (\ii/\Vol)\,\ii^n(-1)^F$: with the factor $\ii^n$ in the lines, the response is $+\ii/\Vol$ times its diagrams, where the imaginary-time rule has $-1/\Vol$. For the bubble this is $(\ii/\Vol)(-1)\sum_{\vect{k}s}\int\dd\omega\,(2\pi)^{-1}\GF_0(\vect{k},\omega)\GF_0(\vect{k}+\vect{q},\omega + \Omega)$, whose frequency integral, done by closing the contour, gives the bubble \cref{eq:4.7-bubble} at $T = 0$ with time-ordered denominators. For a boson the same replacement \cref{eq:D.11-boson} holds, with the time-ordered propagator $\mathsf{D}_0(\vect{Q},t) = -\ii\langle\Phi_0|\mathrm{T}\,\varphi_{\vect{Q},\mathrm{I}}(t)\varphi_{-\vect{Q},\mathrm{I}}(0)|\Phi_0\rangle = -\ii\,\ee^{-\ii\omega_{\vect{Q}}|t|}$, whose transform is $\mathsf{D}_0(\vect{Q},\omega) = 2\omega_{\vect{Q}}/(\omega^2 - \omega_{\vect{Q}}^2 + \ii0^+)$. The imaginary-time rules are the ones used in practice at finite temperature, and the zero-temperature rules are their limit after the frequency sums have been done and continued to the real axis (\cref{sec:4.7}).
